\section{Bac Blanc 2025 — Collège F.-X. Vogt}

\begin{synthese}[Une vraie annale]
Voici, transcrite fidèlement, l'épreuve de \textbf{Baccalauréat Blanc} (séquence 06)
du \textbf{Collège Mgr François-Xavier Vogt}, année 2024/2025, série \textbf{TI}.
Elle réunit, comme à l'examen, une \textbf{Partie I — Programmation web} (HTML,
JavaScript, PHP, MySQL) et une \textbf{Partie II — Programmation procédurale en C}
(structures, fonctions, récursivité). Barèmes d'origine conservés. Cherche d'abord
seul, puis confronte ta copie au corrigé complet.
\end{synthese}

% =====================================================================
%  EN-TÊTE DU SUJET
% =====================================================================
\begin{center}
\fbox{\begin{minipage}{0.93\linewidth}
\centering\sffamily
{\bfseries COLLÈGE Mgr FRANÇOIS-XAVIER VOGT}\\[1pt]
{\footnotesize Département : Informatique}\\[3pt]
\rule{0.6\linewidth}{0.4pt}\\[3pt]
{\large\bfseries BACCALAURÉAT BLANC — Séquence 06}\\[2pt]
{\bfseries Année scolaire : 2024/2025}\\[3pt]
{\bfseries Épreuve : Algorithmique et Programmation}\\[2pt]
{\bfseries Niveau : T\textsuperscript{les} TI}\\[3pt]
\begin{tabular}{@{}l@{\quad}c@{\quad}l@{}}
Durée : \textbf{03 heures} & $\bullet$ & Coefficient : \textbf{04}
\end{tabular}\\[3pt]
\rule{0.6\linewidth}{0.4pt}\\[2pt]
{\footnotesize\itshape L'épreuve comporte deux parties indépendantes, notées chacune sur 10 points.}
\end{minipage}}
\end{center}

% =====================================================================
%  PARTIE I — PROGRAMMATION WEB
% =====================================================================
\subsection*{Partie I — Programmation web \hfill (10 points)}
\addcontentsline{toc}{subsection}{Vogt 2025 — Partie I : Programmation web}

\noindent\textbf{Contexte.} \emph{« Vous êtes stagiaire au sein d'une entreprise locale de
commerce électronique. L'un des projets consiste à développer une application web pour la
gestion des commandes de produits et sa publication via un hébergeur web dynamique.
L'entreprise opte pour une programmation modulaire et vous êtes chargé du module d'ajout de
nouveaux produits. Ce module consiste à créer deux pages web. Une page permet de saisir les
détails du produit, et ces détails seront stockés dans une base de données nommée
\code{ECOM\_DB}, où l'une des tables est nommée \code{PRODUITS}. »}

\subsubsection*{Exercice 1 — Programmation HTML et JavaScript \hfill (03 pts)}

\begin{enumerate}[label=\textbf{\arabic*.}]
\item \textbf{(0,25 pt)} Définir : \emph{Application web}.
\item \textbf{(1,25 pt)} Écrire le code HTML pour afficher le formulaire (nommé
\code{productForm}). La méthode est \code{POST}, et le fichier de traitement est nommé
\code{save\_product.php}. Les noms des champs du formulaire sont respectivement
\emph{« Nom du produit »}, \emph{« Catégorie »}, \emph{« Prix »}, \emph{« Quantité »} et
\emph{« Enregistrer »}. Vous pouvez utiliser un tableau de 05 lignes et 02 colonnes. Inclure
une fonction JavaScript de validation de formulaire nommée \code{validateForm()}.
\item \textbf{(0,5 pt)} Préciser une autre méthode de récupération des données d'un
formulaire, puis donner un avantage de la méthode \code{POST} sur cette autre méthode.
\item \textbf{(0,75 pt)} On désire récupérer les données issues du formulaire et les
enregistrer en vérifiant juste que tous les champs sont renseignés. Si c'est le cas, le
message \emph{« Produit enregistré avec succès »} est affiché dans une boîte de dialogue ;
ou bien \emph{« Échec d'enregistrement du produit »} dans le cas contraire. Écrire en
JavaScript le code de la fonction \code{validateForm()} permettant d'effectuer cette tâche
au clic du bouton \emph{« Enregistrer »}.
\end{enumerate}

\subsubsection*{Exercice 2 — Programmation PHP et MySQL \hfill (07 pts)}

\noindent On s'intéresse maintenant à la base de données \code{ECOM\_DB} et au moyen d'y accéder.

\begin{enumerate}[label=\textbf{\arabic*.}]
\item \textbf{(0,25 pt)} Définir : \emph{Hébergeur web}.
\item \textbf{(0,25 pt)} Donner la signification de \emph{PHP}.
\item \textbf{(0,5 pt)} Donner deux avantages du PHP par rapport à JavaScript.
\item \textbf{(0,75 pt)} Donner les éléments constitutifs d'un environnement intégré de
développement web.
\item \textbf{(0,25 pt)} Indiquer le côté dans lequel sont exécutés les scripts PHP.
\item \textbf{(1,5 pt)} Écrire un code PHP qui, au clic du bouton \emph{« Envoyer »},
récupère les informations du formulaire et les affiche. On prendra pour nom d'hôte
\code{localhost}, l'utilisateur \code{root} et le mot de passe vide.
\end{enumerate}

\noindent Soit l'extrait de code ci-dessous :

\begin{codephp}
<?php
$result = mysqli_query($conn, "SELECT Product_Name, Category, Price FROM PRODUCTS") or die("ERREUR: " . mysqli_error($conn));
while ($row = mysqli_fetch_array($result)) {
    echo "Nom du produit: " . $row["Product_Name"] . "<br/>";
    echo "Catégorie: " . $row["Category"] . "<br/>";
    echo "Prix: " . $row["Price"] . "<br/>";
}
?>
\end{codephp}

\begin{enumerate}[label=\textbf{\arabic*.}, start=7]
\item \textbf{(1 pt)} Donner le rôle des fonctions ci-dessous :
\code{mysqli\_fetch\_assoc()} ; \code{die()}.
\item \textbf{(0,5 pt)} Présenter la différence entre les fonctions
\code{mysqli\_fetch\_array()} et \code{mysqli\_fetch\_row()}.
\item \textbf{(2 pts)} Écrire un code PHP qui permet de supprimer un produit de la base de
données, à partir du nom saisi dans un formulaire adéquat, puis d'afficher le contenu de la
table modifiée dans un tableau HTML.
\end{enumerate}

% =====================================================================
%  PARTIE II — PROGRAMMATION PROCÉDURALE EN C
% =====================================================================
\subsection*{Partie II — Programmation procédurale en C \hfill (10 points)}
\addcontentsline{toc}{subsection}{Vogt 2025 — Partie II : Programmation en C}

\subsubsection*{Exercice 1 \hfill (5 pts)}

\noindent\emph{« Fadil a été choisi pour écrire un programme C permettant de gérer les
stagiaires caractérisés par leur numéro d'inscription (entier), nom (chaîne), prénom
(chaîne), filière (chaîne) et moyenne (réel). On lui a notamment confié la partie consistant
à implémenter trois fonctions : Enregistrement, Modification et Recherche. »} Aider Fadil à
réaliser les tâches ci-dessous.

\begin{enumerate}[label=\textbf{\arabic*.}]
\item \textbf{(1,5 pt)} Donner l'instruction C permettant de déclarer une structure
d'enregistrement nommée \code{stagiaire} pour gérer les stagiaires avec les caractéristiques
données dans le texte.
\item \textbf{(0,5 pt)} Donner l'instruction C permettant de déclarer un tableau \code{T} de
100 stagiaires.
\item \textbf{(0,5 pt)} Donner deux avantages de l'utilisation des fonctions dans un programme.
\item \textbf{(2,5 pts)} Donner en C le code de définition de la fonction \code{Rechercher}
permettant de rechercher un stagiaire par son numéro d'inscription dans un tableau. La
fonction affiche le nom, le prénom et la filière du stagiaire recherché. Dans le cas où le
stagiaire n'est pas présent dans le tableau, la fonction affiche le message :
\emph{« Aucun stagiaire ne correspond à ce numéro d'inscription »}. La fonction prend comme
paramètres le tableau de stagiaires \code{T}, sa taille \code{n} et le numéro d'inscription à
rechercher \code{num\_insc}. L'en-tête de la fonction est :
\code{void Rechercher(struct stagiaire T[], int n, int num\_insc)}.
\end{enumerate}

\subsubsection*{Exercice 2 \hfill (5 pts)}

\noindent On donne le programme C ci-dessous, qui calcule la factorielle d'un nombre (les
numéros de ligne 1 à 27 sont rappelés en commentaire pour pouvoir y faire référence) :

\begin{codec}
/* 1*/  /* Exemple de fonction recurrente. Calcul de la factorielle d'un nombre */
/* 2*/  #include <stdio.h>
/* 3*/  #include <stdlib.h>
/* 4*/  unsigned int f, x;
/* 5*/  unsigned int factorielle(unsigned int a);
/* 6*/  int main() {
/* 7*/      puts("Entrez une valeur entiere entre 1 et 8: ");
/* 8*/      scanf("%d", &x);
/* 9*/      if (x > 8 || x < 1) {
/*10*/          printf("On a dit entre 1 to 8 !");
/*11*/      }
/*12*/      else {
/*13*/          f = factorielle(x);
/*14*/          printf("Factorielle %u egal %u\n", x, f);
/*15*/      }
/*16*/      exit(EXIT_SUCCESS);
/*17*/  }
/*18*/  unsigned int factorielle(unsigned int a) {
/*19*/      if (a == 1) {
/*20*/          return 1;
/*21*/      }
/*22*/      else {
/*23*/          a *= factorielle(a-1);
/*24*/          return a;
/*25*/      }
/*26*/  }
/*27*/
\end{codec}

\begin{enumerate}[label=\textbf{\arabic*.}]
\item Donner le rôle de :
  \begin{enumerate}[label=\textbf{1.\arabic*}]
  \item \textbf{(0,25 pt)} la ligne 5 (la déclaration \code{unsigned int factorielle(unsigned int a);}) ;
  \item \textbf{(0,25 pt)} la ligne 8 (\code{scanf("\%d", \&x);}).
  \end{enumerate}
\item \textbf{(0,5 pt)} Donner la différence fondamentale entre la variable \code{x} et la
variable \code{a} dans ce code.
\item \textbf{(1,5 pt)} Exécuter pas à pas ce programme lorsque la valeur de \code{x = 6}
(présenter la trace, de préférence dans un tableau).
\item \textbf{(1 pt)} Réécrire les instructions de la structure \code{if/else} (lignes 9 à 13)
en utilisant l'opérateur ternaire (\code{ ? : }).
\item \textbf{(1,5 pt)} Réécrire la version itérative de la fonction \code{factorielle}.
\end{enumerate}

% =====================================================================
%  CORRIGÉ
% =====================================================================
\corrige{de l'annale Vogt 2025}

\subsubsection*{Partie I — Exercice 1 (HTML / JavaScript)}

\textbf{1. Définition — Application web.} Une \emph{application web} est un logiciel
applicatif accessible et exécuté à travers un navigateur web, via le réseau Internet ou un
intranet. Côté serveur, elle s'appuie sur un langage dynamique (PHP, etc.) et, le plus
souvent, sur une base de données ; l'utilisateur n'installe rien sur sa machine.

\medskip
\textbf{2. Formulaire HTML (\code{productForm}) avec tableau 5$\times$2.}

\begin{codehtml}
<!DOCTYPE html>
<html lang="fr">
<head>
  <meta charset="utf-8">
  <title>Ajout d'un produit</title>
  <script src="validation.js"></script>
</head>
<body>
  <h2>Nouveau produit</h2>
  <form name="productForm" method="POST" action="save_product.php"
        onsubmit="return validateForm()">
    <table border="1" cellpadding="6">
      <tr>
        <td>Nom du produit</td>
        <td><input type="text" name="Nom du produit"></td>
      </tr>
      <tr>
        <td>Catégorie</td>
        <td><input type="text" name="Catégorie"></td>
      </tr>
      <tr>
        <td>Prix</td>
        <td><input type="number" name="Prix"></td>
      </tr>
      <tr>
        <td>Quantité</td>
        <td><input type="number" name="Quantité"></td>
      </tr>
      <tr>
        <td colspan="2" align="center">
          <input type="submit" name="Enregistrer" value="Enregistrer">
        </td>
      </tr>
    </table>
  </form>
</body>
</html>
\end{codehtml}

\fig{Tableau de 5 lignes $\times$ 2 colonnes ; la validation est déclenchée par \code{onsubmit}.}

\begin{remarque}
Le formulaire porte le nom \code{productForm}, la méthode est \code{POST} et l'attribut
\code{action} pointe vers \code{save\_product.php}. L'attribut \code{onsubmit="return
validateForm()"} bloque l'envoi si la fonction retourne \code{false}.
\end{remarque}

\medskip
\textbf{3. Autre méthode + avantage de POST.} L'autre méthode de récupération des données
d'un formulaire est la méthode \textbf{GET} (les données circulent dans l'URL, après le
\code{?}). Un avantage de \code{POST} sur \code{GET} : avec \code{POST}, les données sont
transmises dans le \emph{corps} de la requête (elles n'apparaissent pas dans l'URL), ce qui
est plus \textbf{sûr/discret} et permet d'envoyer un \textbf{volume de données plus
important} (GET est limité par la longueur de l'URL).

\medskip
\textbf{4. Fonction \code{validateForm()}.}

\begin{codejs}
function validateForm() {
  // Récupération des champs du formulaire productForm
  var f = document.forms["productForm"];
  var nom       = f["Nom du produit"].value;
  var categorie = f["Catégorie"].value;
  var prix      = f["Prix"].value;
  var quantite  = f["Quantité"].value;

  // Vérifie que tous les champs sont renseignés
  if (nom === "" || categorie === "" || prix === "" || quantite === "") {
    alert("Échec d'enregistrement du produit");
    return false;   // empêche l'envoi du formulaire
  } else {
    alert("Produit enregistré avec succès");
    return true;    // autorise l'envoi vers save_product.php
  }
}
\end{codejs}

\begin{remarque}
\code{return false} empêche la soumission (les champs sont incomplets) ; \code{return true}
laisse le formulaire partir vers \code{save\_product.php}. La fonction est appelée par
l'attribut \code{onsubmit="return validateForm()"} placé sur la balise \code{<form>}.
\end{remarque}

\subsubsection*{Partie I — Exercice 2 (PHP / MySQL)}

\textbf{1. Définition — Hébergeur web.} Un \emph{hébergeur web} est une entreprise (ou un
prestataire) qui met à disposition des serveurs connectés en permanence à Internet pour
stocker les fichiers d'un site web et le rendre accessible 24h/24 aux internautes (souvent
avec base de données, nom de domaine, etc.).

\medskip
\textbf{2. Signification de PHP.} \emph{PHP} signifie \textbf{PHP: Hypertext Preprocessor}
(à l'origine \emph{Personal Home Page}).

\medskip
\textbf{3. Deux avantages du PHP par rapport à JavaScript.}
\begin{itemize}
  \item PHP s'exécute \textbf{côté serveur} : il peut se connecter à une base de données et
  manipuler des fichiers du serveur en toute sécurité (le code source n'est pas visible par
  le client) ;
  \item PHP permet de \textbf{générer dynamiquement du HTML} et de gérer la persistance des
  données (sessions, base de données), alors que JavaScript (côté client) ne peut pas
  accéder directement à la base de données.
\end{itemize}

\medskip
\textbf{4. Éléments constitutifs d'un environnement intégré de développement web.} Un EID web
(de type WampServer/XAMPP) regroupe : un \textbf{serveur web} (Apache), un
\textbf{interpréteur PHP} (le langage serveur), un \textbf{SGBD} (MySQL / MariaDB) et un
outil d'administration de la base (\textbf{phpMyAdmin}). On y ajoute un \textbf{éditeur de
code} et un \textbf{navigateur web}.

\medskip
\textbf{5. Côté d'exécution des scripts PHP.} Les scripts PHP sont exécutés \textbf{du côté
serveur} (server side).

\medskip
\textbf{6. Code PHP de récupération et d'affichage des données du formulaire.}

\begin{codephp}
<?php
// save_product.php — récupère et affiche les données du formulaire
if (isset($_POST["Enregistrer"])) {
    // Connexion à la base ECOM_DB
    $conn = mysqli_connect("localhost", "root", "", "ECOM_DB");
    if (!$conn) {
        die("Connexion echouee: " . mysqli_connect_error());
    }

    // Récupération des champs envoyés par le formulaire (méthode POST)
    $nom       = $_POST["Nom du produit"];
    $categorie = $_POST["Catégorie"];
    $prix      = $_POST["Prix"];
    $quantite  = $_POST["Quantité"];

    // Affichage des informations reçues
    echo "Nom du produit : " . $nom . "<br/>";
    echo "Catégorie : "      . $categorie . "<br/>";
    echo "Prix : "           . $prix . " FCFA<br/>";
    echo "Quantité : "       . $quantite . "<br/>";

    mysqli_close($conn);
}
?>
\end{codephp}

\medskip
\textbf{7. Rôle des fonctions.}
\begin{itemize}
  \item \code{mysqli\_fetch\_assoc()} : récupère la \emph{ligne courante} d'un résultat de
  requête \code{SELECT} sous forme de \textbf{tableau associatif}, dont les clés sont les
  \textbf{noms des colonnes} (ex. \code{\$row["Product\_Name"]}). À chaque appel, on avance
  d'une ligne ; elle renvoie \code{NULL} quand il n'y a plus de lignes.
  \item \code{die()} : \textbf{arrête immédiatement l'exécution} du script PHP (équivalent de
  \code{exit()}) ; on lui passe souvent un message qui sera affiché avant l'arrêt. Elle sert
  à interrompre proprement le script en cas d'erreur (échec de connexion ou de requête).
\end{itemize}

\medskip
\textbf{8. Différence entre \code{mysqli\_fetch\_array()} et \code{mysqli\_fetch\_row()}.}
\code{mysqli\_fetch\_row()} renvoie la ligne sous forme de tableau \textbf{indexé
numériquement} uniquement (accès par \code{\$row[0]}, \code{\$row[1]}\dots). En revanche,
\code{mysqli\_fetch\_array()} renvoie par défaut un tableau \textbf{à la fois indexé
numériquement ET associatif} (accès possible par \code{\$row[0]} \emph{ou} par
\code{\$row["nom\_colonne"]}).

\medskip
\textbf{9. Code PHP de suppression d'un produit + ré-affichage du tableau.}

\begin{codephp}
<?php
// delete_product.php — supprime un produit puis réaffiche la table
$conn = mysqli_connect("localhost", "root", "", "ECOM_DB");
if (!$conn) {
    die("Connexion echouee: " . mysqli_connect_error());
}

// Si le formulaire de suppression a été soumis
if (isset($_POST["supprimer"])) {
    $nom = $_POST["nom_produit"];
    $sql = "DELETE FROM PRODUITS WHERE Product_Name = '" . $nom . "'";
    if (mysqli_query($conn, $sql)) {
        echo "<p>Produit supprimé : " . $nom . "</p>";
    } else {
        echo "<p>Erreur : " . mysqli_error($conn) . "</p>";
    }
}
?>

<!-- Formulaire de saisie du nom à supprimer -->
<form method="POST" action="delete_product.php">
  Nom du produit à supprimer :
  <input type="text" name="nom_produit">
  <input type="submit" name="supprimer" value="Supprimer">
</form>

<!-- Affichage du contenu de la table modifiée dans un tableau HTML -->
<table border="1" cellpadding="6">
  <tr>
    <th>Nom du produit</th>
    <th>Catégorie</th>
    <th>Prix</th>
  </tr>
  <?php
  $result = mysqli_query($conn, "SELECT Product_Name, Category, Price FROM PRODUITS");
  while ($row = mysqli_fetch_assoc($result)) {
      echo "<tr>";
      echo "<td>" . $row["Product_Name"] . "</td>";
      echo "<td>" . $row["Category"] . "</td>";
      echo "<td>" . $row["Price"] . "</td>";
      echo "</tr>";
  }
  mysqli_close($conn);
  ?>
</table>
\end{codephp}

\begin{attention}
En production, on ne concatène jamais directement une saisie utilisateur dans une requête
(risque d'injection SQL). On utiliserait une \textbf{requête préparée}. Ici, le code suit le
niveau attendu à l'examen Terminale TI.
\end{attention}

\subsubsection*{Partie II — Exercice 1 (structures et fonctions en C)}

\textbf{1. Déclaration de la structure \code{stagiaire}.}

\begin{codec}
struct stagiaire {
    int   num_insc;      /* numéro d'inscription (entier)   */
    char  nom[50];       /* nom (chaîne de caractères)      */
    char  prenom[50];    /* prénom (chaîne de caractères)   */
    char  filiere[50];   /* filière (chaîne de caractères)  */
    float moyenne;       /* moyenne (réel)                  */
};
\end{codec}

\medskip
\textbf{2. Déclaration d'un tableau \code{T} de 100 stagiaires.}

\begin{codec}
struct stagiaire T[100];
\end{codec}

\medskip
\textbf{3. Deux avantages des fonctions.}
\begin{itemize}
  \item Elles évitent la \textbf{répétition de code} (réutilisation) : on écrit un traitement
  une seule fois puis on l'appelle autant de fois que nécessaire ;
  \item Elles rendent le programme plus \textbf{lisible, structuré et facile à maintenir/}
  \textbf{déboguer} (découpage en sous-problèmes — programmation modulaire).
\end{itemize}

\medskip
\textbf{4. Définition de la fonction \code{Rechercher}.}

\begin{codec}
#include <stdio.h>

void Rechercher(struct stagiaire T[], int n, int num_insc) {
    int i;
    int trouve = 0;   /* drapeau : 0 = non trouvé, 1 = trouvé */

    for (i = 0; i < n; i++) {
        if (T[i].num_insc == num_insc) {
            printf("Nom    : %s\n", T[i].nom);
            printf("Prénom : %s\n", T[i].prenom);
            printf("Filière: %s\n", T[i].filiere);
            trouve = 1;
            break;   /* numéro d'inscription unique : on s'arrête */
        }
    }

    if (trouve == 0) {
        printf("Aucun stagiaire ne correspond à ce numéro d'inscription\n");
    }
}
\end{codec}

\begin{remarque}
On parcourt le tableau du début à la fin. Dès qu'un \code{num\_insc} correspond, on affiche
le nom, le prénom et la filière, on lève le drapeau \code{trouve} et on sort de la boucle. Si
aucune correspondance n'est trouvée après le parcours, on affiche le message d'absence.
\end{remarque}

\subsubsection*{Partie II — Exercice 2 (factorielle récursive)}

\textbf{1. Rôle des lignes.}
\begin{itemize}
  \item \textbf{Ligne 5} — \code{unsigned int factorielle(unsigned int a);} : c'est le
  \textbf{prototype} (déclaration anticipée) de la fonction \code{factorielle}. Il annonce au
  compilateur, avant \code{main}, le nom de la fonction, son type de retour
  (\code{unsigned int}) et le type de son paramètre, pour qu'elle puisse être appelée dans
  \code{main} (ligne 13) alors qu'elle n'est définie qu'ensuite (ligne 18).
  \item \textbf{Ligne 8} — \code{scanf("\%d", \&x);} : \textbf{lit au clavier} un entier
  saisi par l'utilisateur et le \textbf{range dans la variable \code{x}} (l'opérateur
  \code{\&} fournit l'adresse de \code{x}).
\end{itemize}

\medskip
\textbf{2. Différence fondamentale entre \code{x} et \code{a}.} \code{x} est une variable
\textbf{globale} (déclarée ligne 4, hors de toute fonction : visible dans tout le programme),
tandis que \code{a} est une variable \textbf{locale} — c'est le \textbf{paramètre} de la
fonction \code{factorielle} (elle n'existe que pendant l'exécution de la fonction et reçoit
une copie de la valeur passée à l'appel).

\medskip
\textbf{3. Trace d'exécution pour \code{x = 6}.} L'appel \code{factorielle(6)} déclenche une
suite d'appels récursifs jusqu'au cas de base \code{a == 1}, puis les valeurs « remontent ».

\begin{center}\small
\begin{tabular}{@{}clcl@{}}
\toprule
\textbf{Appel} & \textbf{Test \code{a==1} ?} & \textbf{Calcul} & \textbf{Valeur retournée}\\
\midrule
\code{factorielle(6)} & non & \code{6 * factorielle(5)} & \code{6 * 120 = 720}\\
\code{factorielle(5)} & non & \code{5 * factorielle(4)} & \code{5 * 24  = 120}\\
\code{factorielle(4)} & non & \code{4 * factorielle(3)} & \code{4 * 6   = 24}\\
\code{factorielle(3)} & non & \code{3 * factorielle(2)} & \code{3 * 2   = 6}\\
\code{factorielle(2)} & non & \code{2 * factorielle(1)} & \code{2 * 1   = 2}\\
\code{factorielle(1)} & oui & cas de base & \code{1}\\
\bottomrule
\end{tabular}
\end{center}

\noindent En remontant : $1 \to 2 \to 6 \to 24 \to 120 \to 720$. Le programme affiche donc :
\code{Factorielle 6 egal 720}.

\medskip
\textbf{4. Réécriture des lignes 9 à 13 avec l'opérateur ternaire.}

\begin{codec}
(x > 8 || x < 1)
    ? printf("On a dit entre 1 to 8 !")
    : printf("Factorielle %u egal %u\n", x, f = factorielle(x));
\end{codec}

\begin{remarque}
L'opérateur ternaire \code{condition ? valeur\_si\_vrai : valeur\_si\_faux} remplace le
\code{if/else}. Si \code{(x > 8 || x < 1)} est vrai, on affiche le message d'erreur ; sinon
on calcule \code{f = factorielle(x)} et on l'affiche.
\end{remarque}

\medskip
\textbf{5. Version itérative de la fonction \code{factorielle}.}

\begin{codec}
unsigned int factorielle(unsigned int a) {
    unsigned int resultat = 1;
    unsigned int i;
    for (i = 2; i <= a; i++) {
        resultat *= i;   /* resultat = resultat * i */
    }
    return resultat;     /* factorielle(1) et factorielle(0) renvoient 1 */
}
\end{codec}

\begin{remarque}
La boucle multiplie \code{resultat} par tous les entiers de 2 à \code{a}. Pour \code{a = 6} :
$1\times2\times3\times4\times5\times6 = 720$, comme la version récursive.
\end{remarque}

\begin{aretenir}
Au Bac TI, la \textbf{Partie web} attend du code \emph{exact} : un \code{<form>} (méthode,
\code{action}), une validation JavaScript (\code{return false} bloque l'envoi),
\code{mysqli\_connect} + \code{\$\_POST} côté serveur, et l'affichage d'un \code{SELECT} dans
un \code{<table>}. La \textbf{Partie C} attend une \code{struct} bien typée, un parcours de
tableau avec drapeau pour la recherche, et la maîtrise de la \textbf{récursivité} (cas de
base + cas récursif) face à sa version \textbf{itérative}.
\end{aretenir}
